Trust in computation done by third-party software is increasingly
becoming a concern in today's world.  As an example, consider a
software implementing an algorithm for which we already have a
rigorous proof of correctness.  Even when the software is implemented
by trusted and expert programmers, one cannot be sure that the
implementation will never produce an erroneous result.  


Insidious bugs may go undetected for several reasons: for
example, the implementation may be too complex to have been tested
comprehensively, or it may use code from untrusted libraries that have
bugs, or the programmer may have inadvertently introduced a bug.
























In such a scenario, we wish to ask: \emph{How can we provide assurance
that an untrusted implementation of an algorithm has produced a
correct result for a given problem instance?}

One way to answer the above question is to formally model the semantics of
every instruction in the implementation in a theorem proving system

and then prove that the composition of these instructions always
produces the correct result.


This is a challenging and skill-intensive exercise for anything but
the simplest of software, but once completed, it gives very
high assurance about the correctness of results produced by the
implementation \emph{for all instances of the problem being solved.}
A relatively less daunting alternative is to have the algorithm (and
its implementation) output for each problem instance, not only the
result but also a \emph{certificate} of correctness of the result.
The certificate must be independently checkable by an auditor, and a
successful audit must yield a proof of correctness of the computed
result \emph{for the given problem instance}.  We call such
algorithms \emph{auditable algorithms}. Note that a successful audit
for one problem instance does not prove bug-freeness of the
implementation; it only shows that the result computed for this
specific problem instance is correct.  A key requirement for effective
use of auditable algorithms is access to an independent trusted
certificate auditor. Fortunately, a certificate auditor is usually
much simpler to design, and its implementation far easier to prove
correct than the implementation of the original algorithm.  Therefore,
auditable algorithms provide a pragmatic solution to the problem of
assuring correctness for untrusted implementations of algorithms.  Not
surprisingly, auditable algorithms have been studied earlier in the
context of different
problems~\cite{mehlhorn-auditable,Heule13,Wetzler14,Gocht22,Ekici17,Barthe06,Necula98,Magron15,Cheung16}.





In this paper, we present the first detailed study of auditable
algorithms for model counting with provable approximation guarantees--
an important problem with diverse applications.  Our study reveals a
nuanced landscape where the complexity of an auditable algorithm can
be ``traded off'' for suitable measures of complexity of certificate
auditing.  This opens up the possibility of treating certificate audit
complexity as a first-class concern in designing auditable
algorithms.

Before we delve into further details, a brief background on model
counting is useful.  Formally, given a system of constraints
$\varphi$, model counting requires us to count the solutions (or
models) of $\varphi$.  This problem is known to be computationally
hard~\cite{Valiant79}, and consequently, 

significant effort has therefore been invested over the
years in designing approximate model counters.

Over the last
decade, riding on spectacular advances in propositional satisfiability
(SAT) solving techniques (see~\cite{sat-technique} for details),
approximate model counting technology has matured to the point where
we can count solutions of propositional constraints with several
thousands of variables within a couple of hours on a standard laptop,
while providing strong approximation guarantees~\cite{approxmc}.  This
has led to a surge in the number of tools from different application
domains that have started using approximate model counters as
black-box
engines~\cite{klebanov2014,teuber2021quantifying,beck2020automating,zinkus2023mcfil}
to solve various problems.  With increasing such usage comes
increasing demands of assurance on the correctness of counts computed
by approximate counters.  Auditable approximate counting offers a
promising approach to solve this problem.

Approximate model counting algorithms come in two different flavours,
viz. deterministic and randomized.  While randomized counting
algorithms with probably approximately correct (PAC)
guarantees~\cite{Valiant84,Valiant84CACM} have been found to scale
better in practice, { it is not possible to provide a meaningful
certificate of correctness for the result computed by a single run of
a randomized algorithm (see~\cite{mehlhorn-auditable} for a nice
exposition on this topic)}.  Therefore, we focus on deterministic
algorithms for approximate model counting.  Specifically, we ask: (a)
what would serve as a meaningful certificate for the best-known
deterministic approximate counting algorithm~\cite{stockmeyer}?
(b) what is the fine-grained complexity of auditing such a
certificate?  (c) can we re-design the approximate counting algorithm
so as to make certificate auditing more efficient?  and (d) is there a
tradeoff in the complexity of an auditable deterministic approximate
counting algorithm and the fine-grained complexity of certificate
auditing?  Our study shows that there is a nuanced interplay between
the complexity of auditable approximate counting algorithms and that
of auditing the certificates generated by such algorithms. Therefore,
if efficiency of certificate auditing is treated as a first-class
concern, we must re-think the design of auditable
algorithms for approximate counting.  Indeed, we present two new
auditable approximate counting algorithms that
significantly improve the efficiency of certificate auditing at the
cost of making the counting algorithms themselves a bit less
efficient.





The primary contributions of the paper can be summarized as listed
below.  In the following, we wish to count models of a propositional
formula on $n$ variables, upto a constant approximation factor.  




\begin{enumerate}
\item We first show that the best-known deterministic approximate
counting algorithm (from~\cite{stockmeyer}) can be converted to an
auditable algorithm that generates a certificate over $\bigO(n^2\log^2
n)$ variables.  The resulting counting algorithm uses
$\bigO(n\log n)$ invocations of a $\Sigma_2^P$ oracle, as in
Stockmeyer's original algorithm.  An auditor, however,
needs only an invocation of a $\Sigma_2^P$ oracle on formulas
constructed over $\bigO(n^2\log^2 n)$ variables.
\item Next, we develop an deterministic approximate counting algorithm for which the certificate uses only $\bigO(n^2)$ variables.  This reduction is achieved by allowing the counting algorithm access to a more powerful $\Sigma_3^P$ (instead of a
$\Sigma_2^P$) oracle, that is invoked only $\bigO(n)$ times.
Certificate auditing in this case requires a single invocation of a
$\Sigma_2^P$ oracle on a formula constructed over $\bigO(n^2)$
variables.
\item Finally, we present another deterministic approximate
counting algorithm that requires a certificate over only $\bigO(n \log
n)$ variables.  This makes certificate auditing significantly more
efficient compared to the cases above, when viewed through the lens of
fine-grained complexity.  This improvement is obtained by allowing the
counting algorithm to invoke a $\Sigma_3^P$ oracle $\bigO(n \log n)$
times. Certificate auditing now requires one invocation of a
$\Sigma_2^P$ oracle on formulas over $\bigO(n \log n)$ variables.
\end{enumerate}
Our results show that the fine-grained complexity of certificate
checking can be ``traded off'' significantly against the power of
oracles accessible to a determistic approximate counting algorithm.
This opens up the possibility of designing new counting algorithms
that are cognizant of certificate checking complexity. Our
proofs involve reasoning about specially constructed formulas with
qantifier alternations, with incomplete information about parts of the
formula.  
This is quite challenging in general, and one of our technical novelties lies in our application of the probabilistic method in the proofs.



The remainder of the paper is organized as follows.  We present
some preliminaries and notation in Section~\ref{sec:prelim}, and
formalize the notion of an audit for approximate counting
algorithms.  In Section~\ref{sec:stock}, we discuss how Stockmeyer's algorithm can be turned into an auditable algorithm, and also present the corresponding audit algorithm. Sections~\ref{sec:bgp} and \ref{sec:combined} discuss the
design of two new auditable algorithms for approximate counting
with successively small certificate sizes, and the complexity
of auditing these certificates.  Finally, we conclude with
a discussion of our findings in Section~\ref{sec:conclusion}. \shortversion{Proof details, omitted due to space constraints, can be found in a full version of the paper available at~\cite{audit-arxiv}.}

